\SourcePage{694}{697}
\section{Elastic collisions of fast electrons with atoms}

Elastic collisions of fast electrons with atoms may be treated by the Born approximation when the incident electron moves much faster than the electrons in the atom.

Because the electron and the atom differ greatly in mass, the atom may be regarded as stationary during the collision. The reference frame in which the centre of inertia is at rest then coincides with the frame in which the atom is at rest. Thus, in formula (126.7), $\mathbf p$ and $\mathbf p'$ are the electron momenta before and after the collision, $m$ is the electron mass, and the angle $\theta$ is the same as the electron deflection angle $\vartheta$. The potential energy $U(r)$ occurring in (126.7), however, must be specified appropriately.

In \S~126 we calculated the matrix element $U_{\mathbf p'\mathbf p}$ of the interaction energy with respect to the free-particle wave functions before and after the collision. For a collision with an atom, the wave functions describing the atom's internal state must also be included. In elastic scattering this state does not change. Consequently, $U_{\mathbf p'\mathbf p}$ must be defined as the matrix element with respect to the electron wave functions $\psi_{\mathbf p}$ and $\psi_{\mathbf p'}$, taken diagonally with respect to the atomic wave function. In other words, $U$ in formula (126.7) is to be understood as the potential energy of interaction between the electron and the atom, averaged over the latter's wave function. It is equal to $e\varphi(r)$, where $\varphi(r)$ is the potential of the field produced at the point $r$ by the mean charge distribution in the atom.

If $\rho(r)$ denotes the charge density in the atom, the potential $\varphi$ obeys Poisson's equation
\[
 \Delta\varphi=-4\pi\rho(r).
\]

\SourcePage{695}{698}
The matrix element $U_{\mathbf p'\mathbf p}$ sought is essentially the Fourier component of $U$ (and hence of $\varphi$) corresponding to the wave vector $\mathbf q=\mathbf k'-\mathbf k$. Applying Poisson's equation separately to each Fourier component gives
\[
 \Delta(\varphi_{\mathbf q}e^{i\mathbf{qr}})=-q^2\varphi_{\mathbf q}e^{i\mathbf{qr}}
 =-4\pi\rho_{\mathbf q}e^{i\mathbf{qr}},
\]
whence $\varphi_{\mathbf q}=4\pi\rho_{\mathbf q}/q^2$, or
\begin{equation}
 \int\varphi e^{-i\mathbf{qr}}dV=\frac{4\pi}{q^2}\int\rho e^{-i\mathbf{qr}}dV.
 \tag{139.1}
\end{equation}
The charge density $\rho(r)$ consists of the electronic charges and the nuclear charge:
\[
 \rho=-en(r)+Ze\delta(r),
\]
where $en(r)$ is the electron-charge density in the atom. Multiplication by $e^{-i\mathbf{qr}}$ followed by integration yields
\[
 \int\rho e^{-i\mathbf{qr}}dV=-e\int ne^{-i\mathbf{qr}}dV+Ze.
\]
We therefore obtain the following expression for the integral of interest:
\begin{equation}
 \int Ue^{-i\mathbf{qr}}dV=\frac{4\pi e^2}{q^2}[Z-F(q)],
 \tag{139.2}
\end{equation}
where
\begin{equation}
 F(q)=\int ne^{-i\mathbf{qr}}dV
 \tag{139.3}
\end{equation}
defines the quantity $F(q)$, called the \emph{atomic form factor}. It depends both on the scattering angle and on the speed of the incident electron.

Finally, substitution of (139.2) into (126.7) gives the cross-section for elastic scattering of fast electrons by an atom in the form\footnote{We neglect exchange effects between the scattered fast electron and the atomic electrons; that is, we do not symmetrize the wave function of the system. The validity of this neglect is evident beforehand: in the ``exchange integral'', interference between the rapidly oscillating free-particle wave function and the atomic-electron wave function makes the associated contribution to the scattering amplitude small.}
\begin{equation}
 d\sigma=\frac{4m^2e^4}{\hbar^4q^4}[Z-F(q)]^2d o,\qquad
 q=\frac{2mv}{\hbar}\sin\frac\vartheta2.
 \tag{139.4}
\end{equation}
Consider first the limiting case $qa_0\ll1$, where $a_0$ gives the order of magnitude of the atomic dimensions. Small $q$ correspond to small scattering angles,

\SourcePage{696}{699}
$\vartheta\ll v_0/v$, where $v_0\sim\hbar/ma_0$ is of the order of the speeds of the atomic electrons.

Expand $F(q)$ in powers of $q$. The zeroth-order term is $\int n\,dV$, the total number $Z$ of electrons in the atom. The first-order term is proportional to $\int\mathbf r n(r)dV$, and therefore to the mean value of the atomic dipole moment; this value vanishes identically (see \S~75). The expansion must consequently be continued through the second-order term, giving
\[
 Z-F(q)=\frac{q^2}{6}\int nr^2dV;
\]
on substituting this into (139.4), we find
\begin{equation}
 d\sigma=\left|\frac{me^2}{3\hbar^2}\int nr^2dV\right|^2d o.
 \tag{139.5}
\end{equation}
Thus, in the small-angle region the cross section is independent of the scattering angle and is fixed by the mean square of the distances between the atomic electrons and the nucleus.

In the opposite limiting case of large $q$ ($qa_0\gg1$, $\vartheta\gg v_0/v$), the factor $e^{-i\mathbf{qr}}$ in the integrand of (139.3) oscillates rapidly, so that the integral as a whole is nearly zero. Hence $F(q)$ may be neglected in comparison with $Z$, leaving

\begin{equation}
 d\sigma=\left(\frac{Ze^2}{2mv^2}\right)^2\frac{d o}{\sin^4(\vartheta/2)},
 \tag{139.6}
\end{equation}
which is the Rutherford cross-section for scattering by the atomic nucleus.

Let us also calculate the transport cross-section
\begin{equation}
 \sigma_{\mathrm{tr}}=\int(1-\cos\vartheta)d\sigma.
 \tag{139.7}
\end{equation}
For angles $\vartheta\ll v_0/v$, equation (139.5) gives
\[
 d\sigma=\mathrm{const}\cdot\sin\vartheta\,d\vartheta
 \approx\mathrm{const}\cdot\vartheta\,d\vartheta,
\]
where const is independent of $\vartheta$. In this range the integrand of (139.7) is therefore proportional to $\vartheta^3d\vartheta$, and the integral rapidly converges at its lower limit. In the range $1\gg\vartheta\gg v_0/v$, on the other hand,
\[
 d\sigma\approx\mathrm{const}\cdot(d\vartheta/\vartheta^3),
\]
and the integrand is proportional to $d\vartheta/\vartheta$; the integral (139.7) is thus logarithmically divergent.

\SourcePage{697}{700}
It follows that precisely this angular range supplies the main part of the integral, so that integration may be restricted to it. The lower limit must be of the order of $v_0/v$; write it as $e^2/(\gamma\hbar v)$, where $\gamma$ is a dimensionless constant. The result is
\begin{equation}
 \sigma_{\mathrm{tr}}=4\pi\left(\frac{Ze^2}{mv^2}\right)^2
 \ln\frac{\gamma\hbar v}{e^2}.
 \tag{139.8}
\end{equation}
An exact calculation of the constant $\gamma$ requires consideration of scattering through angles $\vartheta>v_0/v$ and cannot be carried out in general. The value of $\sigma_{\mathrm{tr}}$ depends only weakly on this constant, since $\gamma$ occurs inside a logarithm multiplied by the large quantity $\hbar v/e^2$.

For a numerical determination of the form factor of heavy atoms, the Thomas--Fermi distribution of the density $n(r)$ may be used. We have seen that, in the Thomas--Fermi model, $n(r)$ has the form
\[
 n(r)=Z^2f(rZ^{1/3}/b)
\]
(all quantities in this and the following formulas are measured in atomic units). It is readily seen that the integral (139.3), evaluated with this $n(r)$, contains $q$ only in a particular combination with $Z$:
\begin{equation}
 F(q)=Z\varphi(bqZ^{-1/3}).
 \tag{139.9}
\end{equation}
Table~11 gives values of the function $\varphi(x)$ common to all atoms\footnote{It must be kept in mind that this formula does not apply at small $q$, consistently with the fact that the integral of $nr^2$ cannot in practice be evaluated by the Thomas--Fermi method (see the note on p.~563). One should also remember that the Thomas--Fermi model does not reproduce the individual atomic properties that disturb their systematic variation with atomic number.}.


\setcounter{table}{10}
\begin{table}[htbp]
\centering
\caption{Thomas--Fermi atomic factor}
\begin{tabular}{@{}rr|rr|rr@{}}
\toprule
$x$ & $\varphi(x)$ & $x$ & $\varphi(x)$ & $x$ & $\varphi(x)$\\
\midrule
0    & 1.000 & 1.08 & 0.422 & 2.17 & 0.224\\
0.15 & 0.922 & 1.24 & 0.378 & 2.32 & 0.205\\
0.31 & 0.796 & 1.39 & 0.342 & 2.48 & 0.189\\
0.46 & 0.684 & 1.55 & 0.309 & 2.64 & 0.175\\
0.62 & 0.589 & 1.70 & 0.284 & 2.79 & 0.167\\
0.77 & 0.522 & 1.86 & 0.264 & 2.94 & 0.156\\
0.93 & 0.469 & 2.02 & 0.240 &      &      \\
\bottomrule
\end{tabular}
\end{table}

With the atomic form factor (139.9), the cross-section (139.4) becomes
\begin{equation}
 d\sigma=\frac{4Z^2}{q^4}[1-\varphi(bqZ^{-1/3})]^2d o
 =Z^{2/3}\Phi\left(Z^{-1/3}v\sin\frac\vartheta2\right)d o,
 \tag{139.10}
\end{equation}

\SourcePage{698}{701}
The total cross-section can be obtained by integration. The small-$\vartheta$ range makes the chief contribution to the integral. We may therefore write
\[
 d\sigma\approx Z^{2/3}\Phi(Z^{-1/3}v\vartheta/2)\cdot2\pi\vartheta\,d\vartheta,
\]
and extend the integration over $d\vartheta$ to infinity:
\[
 \sigma=2\pi Z^{2/3}\int_0^\infty\Phi\left(Z^{-1/3}\frac{v\vartheta}{2}\right)
 \vartheta\,d\vartheta
 =\frac{8\pi}{v^2}Z^{4/3}\int_0^\infty x\Phi(x)dx.
\]
Thus $\sigma$ has the form
\begin{equation}
 \sigma=\mathrm{const}\cdot\frac{Z^{4/3}}{v^2}.
 \tag{139.11}
\end{equation}
It is similarly easy to verify that the constant $\gamma$ in (139.8) is proportional to $Z^{-1/3}$.

\ProblemHeading

Calculate the cross-section for elastic scattering of fast electrons by a hydrogen atom in its ground state.

\SolutionHeading The ground-state wave function of the hydrogen atom is $\psi=e^{-r}/\sqrt\pi$ (in atomic units), and hence $n=e^{-2r}/\pi$. The angular integration in (139.3), carried out as in the derivation of (126.12), gives

\[
 F=\frac{4\pi}{q}\int_0^\infty n(r)\sin qr\,dr
 =\left(1+\frac{q^2}{4}\right)^{-2}.
\]
Substitution in (139.4) yields
\[
 d\sigma=\frac{4(8+q^2)^2}{(4+q^2)^4}d o,
\]
where $q=2v\sin(\vartheta/2)$. To evaluate the total cross-section, write
\[
 d o=2\pi\sin\vartheta\,d\vartheta=\frac{2\pi}{v^2}q\,dq
\]
and integrate with respect to $q$ from $0$ to $2v$. But $v$ is assumed large and the integral converges, so its upper limit may be extended to infinity. We thus obtain
\[
 \sigma=\frac{7\pi}{3v^2}.
\]
The transport cross section is evaluated as the integral

\[
 \sigma_{\mathrm{tr}}=\frac1{2v^2}\int q^2d\sigma.
\]
On changing the integration variable according to $4+q^2=u$, and replacing the upper limit by infinity everywhere except in the term $du/u$, we obtain
\[
 \sigma_{\mathrm{tr}}=\frac{4\pi}{v^2}\left(\ln v+\frac1{12}\right)
\]
in agreement with (139.8).
